% ============================================================================
%  preamble/engine.tex  --  make one source file work on all three engines
% ----------------------------------------------------------------------------
%  LaTeX has three modern compilers and they handle fonts completely
%  differently:
%
%    pdfLaTeX   old, fast, uses pre-built Type1 fonts. Cannot read .ttf/.otf
%               files, so it CANNOT use Sarabun.
%    XeLaTeX    can use any .ttf/.otf font on disk. Needed for Sarabun.
%    LuaLaTeX   same font abilities as XeLaTeX, plus scripting.
%
%  This file detects which one is running and sets up the text encoding
%  accordingly, so the same main.tex compiles under any of them with no edits.
% ============================================================================

\usepackage{iftex}

\ifPDFTeX
  % ------------------------------------------------------------------------
  %  pdfLaTeX branch
  %  T1 fontenc gives proper accented characters and correct hyphenation.
  %
  %  NOTE: \usepackage[utf8]{inputenc} is deliberately NOT loaded.
  %  UTF-8 has been LaTeX's default input encoding since the 2018 release,
  %  so that line is a no-op on any current TeX Live -- and it is an error
  %  under XeLaTeX/LuaLaTeX. Leaving it out keeps this file engine-neutral.
  % ------------------------------------------------------------------------
  \usepackage[T1]{fontenc}
\else
  % ------------------------------------------------------------------------
  %  XeLaTeX / LuaLaTeX branch
  %  fontspec is what lets LaTeX open .ttf/.otf files directly.
  % ------------------------------------------------------------------------
  \usepackage{fontspec}
  \defaultfontfeatures{Ligatures=TeX}
\fi

% ---------------------------------------------------------------------------
%  TYPEFACES
%  Fira Sans ships with TeX Live in BOTH Type1 and OpenType form, and these
%  two packages pick the right one for whichever engine is running. That is
%  why the same two lines work everywhere and do not need to be inside the
%  branch above.
% ---------------------------------------------------------------------------
\usepackage[sfdefault]{FiraSans}   % body + headings
\usepackage{FiraMono}              % code; sets \ttdefault by itself

\renewcommand{\familydefault}{\sfdefault}

% ---------------------------------------------------------------------------
%  MATHS
%  Fira has no matching maths font, so maths stays in Computer Modern. That
%  is deliberate: it is what most maths papers look like and it reads well
%  on a projector.
% ---------------------------------------------------------------------------
